% Recovery and articulation of the source owners documented in
% PROCEDENCIA_OBSERVADOR_RESPUESTA.md. Originals preserved in full.
% Integration antecedents: APP--TRIT--TPK kernel, realization of
% boundary transport, areal/volumetric readers, and energy clock.
% The representation conditions are stated before each use.

\section{Observer, memory, and relational response}
\label{vii:sec:observador-respuesta}

The reading operation acts on the transport expressed by the
enriched APP--TRIT--TPK state. Its domain retains route, sheet, orientation,
carry, and memory; the result accessible to the observer is a map
from that domain. This distinction makes it possible to describe three
operations jointly: the choice of an oriented representative of the transport,
the inscription of its register in an apparatus, and the local response induced
by the global boundary. Each operation conveys different information.
The following construction specifies their maps and the conditions
under which they can be composed.

\subsection{Boundary transport and intrinsic readout}
\label{vii:sec:observador-ciclo}

The finite representation of the transport provides an oriented cycle
\(C_n\), with \(n\geq1\), vertices \(0,\ldots,n-1\), and edges
\(e_j:j\longrightarrow j+1\pmod n\). We fix the lattice realization
\(L\) of the register, its vector space
\(W=L\otimes_{\mathbb Z}\mathbb R\), a positive inner product, and
the nonzero primitive representative \(v\in L\) of the distinguished
residual class, in the declared chamber. The sum and product sheets of APP
provide the polarizations of the transport; the orientation of the TRIT and the
TPK update precede this linear representation. The cycle,
the lattice realization, and the residual class are the data that this
section receives from the transport constructed above.

A \(0\)-cochain is a function on the vertices with values in
\(W\); a \(1\)-cochain assigns a vector to each oriented edge.
The coboundary acts by
\[
 (d\Phi)(e_j)=\Phi(j+1)-\Phi(j).
\]
We denote by \(\mathbf1_E\) the scalar cochain that is constant on the
edges. For a transport \(\omega\in C^1(C_n;W)\), its accumulated
representation \(\Gamma_\omega\) is obtained by partial sums and
affine interpolation; its closure defect is
\begin{equation}
 B(\omega)=\Gamma_\omega(1)-\Gamma_\omega(0)
 =\sum_{j=0}^{n-1}\omega(e_j).
 \label{vii:eq:observador-defecto}
\end{equation}

\begin{theorem}[Exact decomposition and harmonic mode of the cycle]
\label{vii:thm:observador-descomposicion}
Every transport has a unique decomposition
\[
 \omega=d\Phi+u\mathbf1_E,\qquad \Phi(0)=0,
 \qquad u=\frac1n\sum_j\omega(e_j).
\]
In particular, \(B(d\Phi)=0\) and \(B(\omega)=nu\).
\end{theorem}

\begin{proof}
Let \(y_j=\omega(e_j)-u\). The sum of the \(y_j\) is zero.
Define \(\Phi(k)=\sum_{j=0}^{k-1}y_j\); the value for \(k=n\)
coincides with \(\Phi(0)=0\), so the function is well defined
on the cycle, including its final edge. We have \(d\Phi(e_j)=y_j\).
If \(d\Phi+u\mathbf1_E=d\Psi+w\mathbf1_E\), summing over edges gives
\(nu=nw\). Then \(d(\Phi-\Psi)=0\); the connectivity of the cycle
and the normalization at the initial vertex imply \(\Phi=\Psi\).
The identity \(\sum_j(\Phi(j+1)-\Phi(j))=0\) directly proves
the telescopic cancellation.
\end{proof}

\begin{definition}[Admissible transport and normalized readout]
\label{vii:def:observador-admisible}
The transport is admissible when its harmonic mode belongs to
\(\mathbb Rv\). On the space \(\mathcal T_{\rm adm}\) of such
transports we define the global readout and its density per phase:
\begin{equation}
 \mathfrak a_v(\omega)
 =\frac{\langle B(\omega),v\rangle}{\langle v,v\rangle},
 \qquad
 \mu_v(\omega)=\frac{\mathfrak a_v(\omega)}n.
 \label{vii:eq:observador-lectura}
\end{equation}
We have \(B(\omega)=\mathfrak a_v(\omega)v\).
\end{definition}

\begin{theorem}[Determination of the intrinsic readout]
\label{vii:thm:observador-unicidad}
The quotient
\(\mathcal T_{\rm adm}/dC^0(C_n;W)\) is one-dimensional, generated
by the class of \(v\mathbf1_E\). The map \(\mu_v\) is the
unique linear functional on \(\mathcal T_{\rm adm}\) satisfying
\[
 F(\omega+d\Phi)=F(\omega),\qquad F(v\mathbf1_E)=1.
\]
\end{theorem}

\begin{proof}
The preceding decomposition writes each admissible transport as
\(d\Phi+\lambda v\mathbf1_E\). Its class modulo exact cochains
is determined by \(\lambda\). The class of \(v\mathbf1_E\)
is nonzero because its sum is \(nv\neq0\). By linearity and
invariance, every functional with the stated properties has value
\(F(\omega)=\lambda F(v\mathbf1_E)=\lambda\). Formula
\eqref{vii:eq:observador-lectura} gives precisely \(\mu_v=\lambda\).
The global readout is \(\mathfrak a_v=n\mu_v\).
\end{proof}

\begin{definition}[Boundary observer]
\label{vii:def:observador-borde}
With the polarization of the expressed transport fixed, a boundary observer
is a pair \(\mathcal O=(\varepsilon,\Psi)\), where
\(\varepsilon\in\{+1,-1\}\) and \(\Psi\in C^0(C_n;W)\).
Its reading action is
\begin{equation}
 \omega^{\mathcal O}=\varepsilon\omega+d\Psi.
 \label{vii:eq:observador-accion}
\end{equation}
The sign fixes the orientation and the cochain fixes the trivialization.
Two choices of polarization are compared through their expressed
transport; relation \eqref{vii:eq:observador-accion} describes
exactly those representatives that differ by orientation and coboundary.
\end{definition}

\begin{proposition}[Covariance of the readout]
\label{vii:prop:observador-covarianza}
The action preserves admissibility and satisfies
\[
 B(\omega^{\mathcal O})=\varepsilon B(\omega),\quad
 \mathfrak a_v(\omega^{\mathcal O})=\varepsilon\mathfrak a_v(\omega),
 \quad \mu_v(\omega^{\mathcal O})=\varepsilon\mu_v(\omega).
\]
Trivializations with the same orientation express the same closure.
\end{proposition}

\begin{proof}
The sum of the coboundary \(d\Psi\) is zero. Therefore, the harmonic mode
is transformed into \(\varepsilon u\), which still belongs to
\(\mathbb Rv\), and the defect is transformed into \(\varepsilon B\).
The two scalar equalities follow from the normalized projection.
\end{proof}

The unit of this readout is the residual representative \(v\), with its
internal normalization. The freedom of trivialization acts on the
cochain representative; the closure information resides in its
class. The description makes it possible to compare observers while keeping
the genealogical transport, its linear representation, and the scalar
that evaluates the residual class separate.

\subsection{Boundary identity and cellular interior}
\label{vii:sec:observador-stokes}

Let \(K\) be a finite oriented two-dimensional cell complex,
equipped with a fundamental \(2\)-chain \(c_K\) such that
\(\partial c_K=[C_n]\). This condition expresses that the interior
incidences cancel with their multiplicities and orientations.
Let \(\Omega\in C^1(K;W)\) be an extension of \(\omega\), and
\(F_\Omega=d\Omega\) its additive discrete curvature.

\begin{theorem}[Cellular holographic identity of the readout]
\label{vii:thm:observador-stokes}
With the preceding data,
\begin{equation}
 B(\omega)=\langle\Omega,\partial c_K\rangle
          =\langle d\Omega,c_K\rangle,
 \qquad
 \mathfrak a_v(\omega)
 =\frac{\langle\langle d\Omega,c_K\rangle,v\rangle}
        {\langle v,v\rangle}.
 \label{vii:eq:observador-stokes}
\end{equation}
In particular, \(F_\Omega=0\) implies \(\mathfrak a_v(\omega)=0\).
If the closure is nonzero, every extension of the transport on that complex
has nonzero discrete curvature.
\end{theorem}

\begin{proof}
The pairing between chains and cochains satisfies
\(\langle d\Omega,c_K\rangle=\langle\Omega,\partial c_K\rangle\).
On expansion, each interior edge appears with the sum of the incidence
numbers of the adjacent cells; the condition on
\(\partial c_K\) leaves exactly the oriented boundary cycle.
The restriction of \(\Omega\) to that cycle is \(\omega\), whose sum
is \(B(\omega)\). Projection onto \(v\) gives the second formula
and the flatness implication.
\end{proof}

The equality compares two evaluations of the same extended cochain:
the accumulated boundary defect and the coboundary integrated over the interior.
Its curvature takes values in \(W\) and uses the additive cellular coboundary.
Realization by a geometric connection additionally preserves
its own coefficient space and transport law.

\subsection{Isometric inscription of the register in the apparatus}
\label{vii:sec:observador-aparato}

Let \(\mathfrak M_n\) be a nonempty finite set of registers
expressed at level \(n\), with result, route, carry, sheet, and memory.
Its register space is
\(\mathcal H_n=\ell^2(\mathfrak M_n)\), with basis \(|m\rangle\).
The visible update \(F_n\) has a register lift
\[
 \widehat F_n(m)=(F_n(m),m).
\]
As registers of the next level we take a set containing these
pairs. The second coordinate retains the antecedent and makes the
lift injective. This construction specifies the memory retained by
the inscription, including when the visible update merges values.

In a finite apparatus space \(\mathcal K_{A,n}\), we fix an
orthonormal collection \((|a_m\rangle)_{m\in\mathfrak M_n}\) and define
\[
 \iota_n:\mathcal H_n\longrightarrow\mathcal K_{A,n},
 \quad \iota_n|m\rangle=|a_m\rangle,
 \qquad
 V_n:\mathcal H_n\longrightarrow\mathcal H_{n+1},
 \quad V_n|m\rangle=|\widehat F_n(m)\rangle.
\]
Let \(\mathcal P_n=\iota_n\mathcal H_n\) be the pointer subspace.

\begin{theorem}[Compatibility between updating and inscription]
\label{vii:thm:observador-entrelazador}
The maps \(\iota_n\) and \(V_n\) are isometries. There is a unique
isometry \(W_n^0:\mathcal P_n\longrightarrow\mathcal K_{A,n+1}\)
satisfying
\begin{equation}
 W_n^0\iota_n=\iota_{n+1}V_n.
 \label{vii:eq:observador-cuadrado}
\end{equation}
If \(\dim\mathcal K_{A,n+1}\geq\dim\mathcal K_{A,n}\),
it extends to an isometry on all of \(\mathcal K_{A,n}\).
Using auxiliary complements that equalize the dimensions gives
a unitary extension between the enlarged spaces.
\end{theorem}

\begin{proof}
The images of the bases under \(\iota_n\) are orthonormal; the
images under \(V_n\) are also orthonormal by the injectivity of
\(\widehat F_n\). For \(z=\iota_n\psi\) define
\(W_n^0z=\iota_{n+1}V_n\psi\). The representation of \(z\) is unique,
and the three maps preserve its norm. The formula determines the map on
a basis of \(\mathcal P_n\) and proves the square.
If the target dimension is sufficient, the orthogonal complement of
\(W_n^0\mathcal P_n\) contains an orthonormal collection of the size of
\(\mathcal P_n^\perp\); sending a basis of the latter to that collection
constructs the extension. Finally, both spaces are enlarged to
the same dimension and the two prescribed collections are completed to
orthonormal bases. The map between those bases is unitary.
\end{proof}

When the register includes labels of a finite instrument, this same
inscription realizes its environment: if the operators
\(K_{y,a}:\mathcal H_S\to\mathcal H'_S\) satisfy
\(\sum_{y,a}K_{y,a}^*K_{y,a}=I\), the map
\[
 V_{\rm inst}\psi=\sum_{y,a}K_{y,a}\psi\otimes|y,a\rangle
\]
is isometric, since its squared norm is
\(\sum_{y,a}\langle\psi,K_{y,a}^*K_{y,a}\psi\rangle=\|\psi\|^2\).
The inscription \(\iota\) transports each label to its corresponding
pointer. The model thus fixes the premeasurement register; a
material device additionally requires its interaction and calibration.

\subsection{Outcome map and record conditioning}
\label{vii:sec:salida-condicionamiento}

The preceding inscription composes with the apparatus readout on the
complete genealogical state. Coupling retains the operation that produces
the correlation; the outcome map expresses the result; conditioning
restricts the extensions compatible with that record.

Let \((\Omega_S,\Sigma_S)\) and \((\Omega_M,\Sigma_M)\) be the
measurable spaces of genealogical states of the system and apparatus,
obtained from their compatible cylinders and records. A measurement
context is written \(\mathfrak M_a=(U_a,q_a,\mathcal C_a,m_0)\):
\(m_0\) is the apparatus preparation, \(\mathcal C_a\) specifies
compatibilities and the extension horizon, and
\[
 U_a:\Omega_S\times\Omega_M\longrightarrow\Omega_S\times\Omega_M,
 \qquad q_a:\Omega_M\longrightarrow Y_a
\]
are, respectively, the measurable coupling and the measurable readout into
the outcome space \((Y_a,\Sigma_{Y_a})\). In a reversible regime,
\(U_a\) is bijective and bimeasurable; the corresponding unitarity is
also required in the Hilbert-space realization. The individual outcome is the
composition
\begin{equation}
 \operatorname{Out}_a(x)
 =q_a\bigl(\operatorname{pr}_M U_a(x,m_0)\bigr).
 \label{vii:eq:out-genealogico}
\end{equation}
It is determined by the complete state, apparatus preparation, and
context. The statistical description of a preparation is another readout
of the set of compatible states.

\begin{proposition}[Outcome measure and restriction to a fibre]
\label{vii:prop:medida-salida}
Let \(\Omega_a\in\Sigma_S\) be the set of states compatible with the
preparation and context, endowed with a probability measure \(\mu_a\)
on \((\Omega_a,\Sigma_S|_{\Omega_a})\).
The map \(\operatorname{Out}_a|_{\Omega_a}\) is measurable and induces
the outcome probability measure
\begin{equation}
 \nu_a(B)=\mu_a\bigl(\operatorname{Out}_a^{-1}(B)\cap\Omega_a\bigr),
 \qquad B\in\Sigma_{Y_a}.
 \label{vii:eq:medida-imagen-salida}
\end{equation}
For an outcome \(y\) whose singleton is measurable, the fibre
\(\Omega_{a,y}=\{x\in\Omega_a:\operatorname{Out}_a(x)=y\}\)
is measurable. Whenever \(\mu_a(\Omega_{a,y})>0\),
\begin{equation}
 \mu_{a,y}(E)=
 \frac{\mu_a(E\cap\Omega_{a,y})}{\mu_a(\Omega_{a,y})}
 \label{vii:eq:condicionamiento-salida}
\end{equation}
for \(E\in\Sigma_S|_{\Omega_a}\) defines a probability measure
concentrated on that fibre.
\end{proposition}
\begin{proof}
The inclusion \(x\mapsto(x,m_0)\), coupling, apparatus projection, and
readout are measurable; their composition is therefore measurable.
Preimages preserve disjoint unions and the whole space, so
\eqref{vii:eq:medida-imagen-salida} is countably additive with total mass
one. The fibre is the preimage of the singleton \(\{y\}\) within
\(\Omega_a\). Restricting \(\mu_a\) to that fibre gives a positive,
countably additive measure; division by its positive mass normalizes it.
\end{proof}

For a prefix \(s\) of depth \(n\), the extensions following the record
form
\[
 \operatorname{Fut}_{a,y}(s)
 =\{x\in\Omega_{a,y}:x|_n=s\}.
\]
Restriction preserves the prefix already traversed and the coupling
record, while changing the collection of admissible extensions. Outcome
\eqref{vii:eq:out-genealogico} acts on complete states; the outcome law
\eqref{vii:eq:medida-imagen-salida} is obtained by pushing the preparation
measure forward through that map. A density-matrix representation expresses the statistical level of the
quantum realization and retains its conditions of correspondence with
genealogical fibres.


\subsection{Accessible algebra and pointer complement}
\label{vii:sec:observador-expectativa}

Define \(P_m=|a_m\rangle\langle a_m|\),
\(P=\sum_mP_m\), and \(P_\perp=I-P\). The algebra of the pointer
subspace is \(P\mathcal B(\mathcal K_{A,n})P\), whose unit is \(P\).
On it, the diagonal reader is
\[
 \mathcal D_P(A)=\sum_mP_mAP_m.
\]

\begin{proposition}[Register expectation and unital extension]
\label{vii:prop:observador-pinzamiento}
The map \(\mathcal D_P\) is a completely positive,
trace-preserving, idempotent conditional expectation on the
pointer algebra, with \(\mathcal D_P(P)=P\). Its image is
\(\bigoplus_m\mathbb CP_m\).
On the algebra of the full apparatus, the map
\begin{equation}
 \mathcal E_P(A)=\sum_mP_mAP_m+P_\perp AP_\perp
 \label{vii:eq:observador-complemento}
\end{equation}
is a unital, completely positive, trace-preserving
conditional expectation. Its image is
\[
 \left(\bigoplus_m\mathbb CP_m\right)
 \oplus\mathcal B(\mathcal P_n^\perp).
\]
Both maps are orthogonal projections for the
Hilbert--Schmidt inner product. Their spectra are contained in \(\{0,1\}\).
When eliminated components exist, the separation between the
fixed subspace and the kernel is one.
\end{proposition}

\begin{proof}
The projectors \(P_m\) are mutually orthogonal. In the full
apparatus, the collection obtained by adding \(P_\perp\) sums to \(I\). Each summand
\(A\mapsto P_mAP_m\) is completely positive: after any
matrix amplification it acts by congruence with \(I\otimes P_m\).
The same holds for \(P_\perp\). The sum retains that property.
The cyclicity of the trace and the sum of the projectors prove
trace preservation on each domain. Orthogonality annihilates the
cross terms upon iteration of the map, so it is idempotent.
Its fixed elements are precisely the stated blocks, and the
map is bimodular with respect to that subalgebra.
Moreover,
\(\operatorname{Tr}(A^*P_mBP_m)
 =\operatorname{Tr}((P_mAP_m)^*B)\),
which gives self-adjointness for Hilbert--Schmidt. A self-adjoint
projection has only the spectral values zero and one when
both subspaces are nonzero. If only the fixed subspace remains,
the map is the identity.
\end{proof}

The sum \(\sum_mP_mAP_m\), applied without the complement to the full
apparatus, sends \(I\) to \(P\). Its natural unit belongs to the
pointer algebra. Formula \eqref{vii:eq:observador-complemento} also retains
the auxiliary block; an orthonormal resolution of that block
provides, when desired, diagonalization of the entire apparatus.

If \(U\) is a unitary transformation of the apparatus and
\(P'_m=UP_mU^*\), then
\[
 \mathcal E_{P'}(UAU^*)=U\mathcal E_P(A)U^*.
\]
The equality follows by substituting each projector in the sum.
Likewise, for \(A\in\mathcal B(\mathcal P_n)\), the square
\eqref{vii:eq:observador-cuadrado} implies
\[
 \mathcal D_{P_{n+1}}(W_n^0A(W_n^0)^*)
 =W_n^0\mathcal D_{P_n}(A)(W_n^0)^*.
\]
Indeed, each matrix unit \(|a_m\rangle\langle a_{m'}|\)
is transported to the unit corresponding to the distinct registers
\(\widehat F_n(m),\widehat F_n(m')\); the diagonal reader retains
exactly the terms with \(m=m'\).

\subsection{Factorization of the response through the effective boundary}
\label{vii:sec:observador-factorizacion}

Let \(\mathcal X_\infty^{\rm enr}\) be the domain of enriched
compatible states, and let
\(b:\mathcal X_\infty^{\rm enr}\longrightarrow B_\infty\) be its global
boundary trace. We denote by \(B_{\rm ef}=\operatorname{im}b\)
the effectively expressed boundary. For a response space
\(\mathcal V_v\), consider
\(I_v:\mathcal X_\infty^{\rm enr}\longrightarrow\mathcal V_v\).

\begin{theorem}[Criterion for relational response]
\label{vii:thm:observador-factorizacion}
There is a unique map \(\overline I_v:B_{\rm ef}\to\mathcal V_v\)
such that \(I_v=\overline I_v\circ b\) if and only if
\begin{equation}
 b(x)=b(x')\quad\Longrightarrow\quad I_v(x)=I_v(x').
 \label{vii:eq:observador-fibra}
\end{equation}
Uniqueness extends to all of \(B_\infty\) when \(b\) is
surjective.
\end{theorem}

\begin{proof}
A composition through \(b\) is constant on each fibre. Conversely,
for \(\beta\in B_{\rm ef}\), define
\(\overline I_v(\beta)=I_v(x)\) with \(b(x)=\beta\).
Condition \eqref{vii:eq:observador-fibra} guarantees independence
of the representative; every value of \(\overline I_v\) is determined
by some state in that fibre. If \(b\) is surjective,
\(B_{\rm ef}=B_\infty\). Otherwise, the composition equation
determines only the values on \(B_{\rm ef}\).
\end{proof}

This criterion is set-theoretic. For a continuous factorization, one
additionally retains the topologies and the quotient-map condition
on \(b:\mathcal X_\infty^{\rm enr}\to B_{\rm ef}\); under this condition,
continuity of \(I_v\) is equivalent to continuity of its factor. This clarification
separates the determination of a readout from its regularity properties.

Let \(\ell_v:\mathcal X_\infty^{\rm enr}\to\mathcal L_v\) be the
restriction accessible to a local neighbourhood. Effective global
dependence is verified by two states satisfying
\begin{equation}
 \ell_v(x)=\ell_v(x'),\qquad I_v(x)\neq I_v(x').
 \label{vii:eq:observador-testigo-global}
\end{equation}
If \(I_v\) factors through \(b\), these two states have distinct
boundaries. Moreover, the response cannot be written as a function of
\(\ell_v\), since that function would assign the same value to both states.
The difference between boundaries thus acquires a verifiable operator effect.

\subsection{Ambivalence weights and self-inertial response}
\label{vii:sec:observador-ambivalencia}

For the general formula we consider positive readers
\(A_\partial,V_\partial\) and an energy reader \(Q\)
defined on the declared effective boundary domain.
We next make explicit a realization on
\(B_{\rm av}\subseteq B_{\rm ef}\). Memory remains as a coordinate
\(\mathsf M_\partial(\beta)\) of the boundary. The ambivalence
hinge provides the two orientations
\[
 r_{\rm av}=\frac{\pi_{\rm HMT}}{\varphi_{\rm HMT}^2},\qquad
 r_{\rm av}^{J}=\frac{\varphi_{\rm HMT}}{\pi_{\rm HMT}^2}.
\]
The superscript \(J\) expresses interchange of the generating pair.
\paragraph{Normalized realization of the boundary readers.}
The APP--TRIT--TPK modal incidence, developed in the regional
construction, provides the path envelope with matrix
\[
 F_{\rm av}=\begin{pmatrix}0&1\\1&1\end{pmatrix}.
\]
The positive vector \(v=(1,\varphi_{\rm HMT})^{\mathsf T}\)
satisfies \(F_{\rm av}v=\varphi_{\rm HMT}v\). Its normalized
transition and stationary weights are
\begin{equation}
 P_{ij}=\frac{(F_{\rm av})_{ij}v_j}{\varphi_{\rm HMT}v_i},
 \qquad
 P=\begin{pmatrix}0&1\\
 \varphi_{\rm HMT}^{-2}&\varphi_{\rm HMT}^{-1}\end{pmatrix},
 \qquad
 (\mu_{\rm ar},\mu_{\rm vol})
 =\frac{(1,\varphi_{\rm HMT}^{2})}{1+\varphi_{\rm HMT}^{2}}.
 \label{vii:eq:lectores-parry}
\end{equation}
The rows of \(P\) sum to one because
\(\varphi^{-2}+\varphi^{-1}=1\); the stationary equation
\(\mu_{\rm ar}=\mu_{\rm vol}\varphi^{-2}\), together with
normalization, determines the stated weights. This Parry measure
weights the paths of the envelope; the chronological frequency
\((1/3,2/3)\) counts the instants of the modal calendar and retains
its distinct function.

The common-memory-amplitude realization is defined on the sector
\(\mathcal X_{\rm av}\subset\mathcal X_\infty^{\rm enr}\) in which
the readers of the two sheets have the factorization
\begin{equation}
 \mathcal L_{\rm vol}(x)=\mu_{\rm vol}\mathcal M(x),\qquad
 \mathcal L_{\rm ar}(x)=\pi_{\rm HMT}\mu_{\rm ar}\mathcal M(x),
 \qquad \mathcal M(x)>0.
 \label{vii:eq:lectores-amplitud-comun}
\end{equation}
Both values belong to the same positive memory-amplitude
line. The areal representation carries the regional mark
\(\pi_{\rm HMT}\) of the return. The dimensionless readout is obtained
by dividing both sections by \(\mathcal L_{\rm vol}(x)\).
In this way, sums of weights are taken between quantities of the
same type; a possible realization as physical areas and volumes
is composed afterwards with their respective dimensional transducers.

\begin{proposition}[Descent of the normalized readers and the energy clock]
\label{vii:prop:lectores-frontera-explicitos}
On \(B_{\rm av}=b(\mathcal X_{\rm av})\), the rules
\begin{equation}
 A_\partial(b(x))
 =\frac{\mathcal L_{\rm ar}(x)}{\mathcal L_{\rm vol}(x)}
 =r_{\rm av},\qquad V_\partial(b(x))=1
 \label{vii:eq:lectores-normalizados}
\end{equation}
define two positive readers independent of the representative.
With an orientation \(a\) of the action section fixed, the decision
energy of \eqref{v:ant:eq:ii-kb-aut-energia} determines on states
\begin{equation}
 q_a(x)=\frac{h_a(x)}{108t_0(x)}
       =\frac{\hbar_a(x)}{t_P(x)}\in\mathcal L_E^+,
 \qquad t_P=\frac{54}{\pi_{\rm HMT}}t_0.
 \label{vii:eq:lector-q-accion}
\end{equation}
We write
\[
 b_{\rm av}:=b|_{\mathcal X_{\rm av}}:
 \mathcal X_{\rm av}\longrightarrow B_{\rm av}.
\]
There is a unique reader \(Q:B_{\rm av}\to\mathcal L_E^+\) such that
\(q_a|_{\mathcal X_{\rm av}}=Q\circ b_{\rm av}\)
if and only if, for any \(x,x'\in\mathcal X_{\rm av}\),
\begin{equation}
 b(x)=b(x')\ \Longrightarrow\
 \frac{h_a(x)}{t_0(x)}=\frac{h_a(x')}{t_0(x')}.
 \label{vii:eq:descenso-reloj}
\end{equation}
In particular, a chart with fixed action and duration satisfies this
condition and expresses \(Q=h_a/(108t_0)\).
\end{proposition}
\begin{proof}
The quotient of \eqref{vii:eq:lectores-amplitud-comun} cancels the
common amplitude and equals
\(\pi_{\rm HMT}\mu_{\rm ar}/\mu_{\rm vol}
=\pi_{\rm HMT}/\varphi_{\rm HMT}^2\).
The result is the same for all representatives of a fibre
of \(b_{\rm av}\) within \(\mathcal X_{\rm av}\), and it is positive.
Dividing both readers by the same positive factor
preserves the weighted quotients in which they occur.
For the clock, the factorization criterion of
Theorem~\ref{vii:thm:observador-factorizacion}, applied to \(q_a\),
is exactly equivalent to \eqref{vii:eq:descenso-reloj}. Fixed sections
satisfy it; so do variable sections whose
quotient descends through \(b_{\rm av}\). Uniqueness follows from the definition
of \(B_{\rm av}\) as the effective image.
\end{proof}

The value of the reader \(Q\) belongs to the energy line
\(\mathcal L_E=\mathcal L_{\rm act}\otimes\mathcal L_T^*\).
With \(E_1,E_2\) in that same line,
\(E_1E_2/Q^2\) is dimensionless and independent of the positive
energy basis. The symbol \(Q\) in this section denotes exclusively
that clock energy; the heat quantity in the Landauer
realization belongs to the corresponding thermal balance.
The common-amplitude factorization specifies a concrete
realization of the general readers. Every other domain retains its
boundary maps and its own verification of descent.
The full memory remains in the enriched history and in its
boundary representation, even though it cancels in the normalized quotient.


Define
\[
 w_A(\beta)=
 \frac{A_\partial(\beta)r_{\rm av}}
 {A_\partial(\beta)r_{\rm av}+V_\partial(\beta)r_{\rm av}^{J}},
 \qquad w_V(\beta)=1-w_A(\beta).
\]
In realization \eqref{vii:eq:lectores-normalizados}, the weights
admit the exact evaluation
\begin{equation}
 w_A=\frac{r_{\rm av}^{\,2}}{r_{\rm av}^{\,2}+r_{\rm av}^{J}}
     =\frac{\pi_{\rm HMT}^{4}}
            {\pi_{\rm HMT}^{4}+\varphi_{\rm HMT}^{5}},\qquad
 w_V=\frac{\varphi_{\rm HMT}^{5}}
            {\pi_{\rm HMT}^{4}+\varphi_{\rm HMT}^{5}}.
 \label{vii:eq:pesos-amplitud-comun}
\end{equation}
The substitution retains both factors of the definition: the normalized
areal reader and its subsequent weighting by \(r_{\rm av}\).
The generation and regional marking of the return are performed once.
Multiplying numerator
and denominator by \(\varphi_{\rm HMT}^4\pi_{\rm HMT}^2\)
gives the final expression. These weights remain constant in the
common-amplitude sector. With fixed projectors and energies, a
variation in the response may arise from the reader \(Q\);
the witness of global dependence additionally requires states with
the same local restriction and different responses as in
\eqref{vii:eq:observador-testigo-global}.


Let \(P_A,P_V\) be complementary orthogonal projections on a
response space \(\mathcal H_v\), and let \(E_1,E_2>0\) be
two internal energies fixed for the comparison. The response is
\begin{equation}
 \overline I_v(\beta)=\frac{E_1E_2}{Q(\beta)^2}
       \bigl(w_A(\beta)P_A+w_V(\beta)P_V\bigr),
 \qquad I_v=\overline I_v\circ b.
 \label{vii:eq:observador-respuesta-ponderada}
\end{equation}
Energies that are explicit readers of
\(\beta\) are also admitted; in that case they are retained as such in the formula.

\begin{proposition}[Positivity and covariance of the response]
\label{vii:prop:observador-respuesta-positiva}
Response \eqref{vii:eq:observador-respuesta-ponderada} is
self-adjoint, positive definite, and constant on the fibres of \(b\).
Its eigenvalues on the nonzero sectors are
\[
 \lambda_A=E_1E_2Q^{-2}w_A,\qquad
 \lambda_V=E_1E_2Q^{-2}w_V.
\]
The full interchange
\((A_\partial,r_{\rm av},P_A)
 \leftrightarrow(V_\partial,r_{\rm av}^{J},P_V)\)
permutes the two summands and preserves the operator. A unitary
transport of the response space conjugates that operator and preserves
its spectrum with multiplicities.
\end{proposition}

\begin{proof}
Positivity of the readers gives \(0<w_A,w_V<1\). For
\(z=z_A+z_V\), with \(z_A=P_Az\), \(z_V=P_Vz\),
\[
 \langle z,\overline I_vz\rangle
 =\frac{E_1E_2}{Q^2}
   (w_A\|z_A\|^2+w_V\|z_V\|^2)>0
\]
if \(z\neq0\). The orthogonal decomposition proves the spectral
expressions. Exclusive dependence on \(\beta\) gives constancy
on fibres and the factorization. Interchanging the triples
simultaneously interchanges numerators, weights, and projectors, thereby
preserving the sum. For a unitary transport \(U\), the substitution
\(P_A\mapsto UP_AU^*\), \(P_V\mapsto UP_VU^*\) produces
\(U\overline I_vU^*\), with the same spectrum.
\end{proof}

An interchange implemented within the same space by a unitary
sending \(P_A\) to \(P_V\) requires both sectors to have
equal dimension. Interchanging the two triples as labels in the
formula is valid with the respective dimensions of each sector.

Condition \eqref{vii:eq:observador-testigo-global} requires comparison of
the responses. For example, simultaneously multiplying
\(A_\partial,V_\partial\) by the same positive factor preserves
both weights; if \(E_1,E_2,Q\) are also preserved, the response
remains identical. By contrast, with equal weights and energies, replacing
\(Q\) by \(qQ\), \(q>0\), multiplies the response by \(q^{-2}\).
If two states with the same local restriction realize that change with
\(q\neq1\), they provide the global witness. These are operator
separation conditions; the existence of the states is verified
in the boundary domain being realized.

\subsection{Connection of the joint state and expressed acceleration}
\label{vii:sec:observador-autoinercia}

In a differentiable realization of the joint state
\(\Gamma=(x,\mathcal O,m)\), let \(N\) be its realization space,
\(\nabla^{\rm tot}\) its connection, and \(\pi_S:N\to N_S\) a
differentiable projection onto the subsystem, equipped with connection
\(\nabla^S\). The connection and projection retain the provenance
of the transports and readers they realize. The observed frame is
obtained through
\[
 \mathfrak F_{\mathcal O}(\Gamma)
 =\operatorname{Res}_{\mathcal O}(\Gamma,\nabla^{\rm tot}).
\]
For a twice-differentiable trajectory \(\Gamma(t)\),
write \(\gamma_S=\pi_S\circ\Gamma\).

\begin{definition}[Self-inertia and projection defect]
\label{vii:def:observador-autoinercia}
A trajectory is self-inertial with respect to the declared connection
when \(\nabla^{\rm tot}_{\dot\Gamma}\dot\Gamma=0\).
The defect of its projection is
\begin{equation}
 \mathcal K_{\pi_S}(\dot\Gamma,\dot\Gamma)
 =\nabla^S_{\dot\gamma_S}\dot\gamma_S
  -d\pi_S\bigl(\nabla^{\rm tot}_{\dot\Gamma}\dot\Gamma\bigr).
 \label{vii:eq:observador-defecto-proyeccion}
\end{equation}
\end{definition}

\begin{proposition}[Subsystem acceleration]
\label{vii:prop:observador-aceleracion}
For a self-inertial trajectory,
\[
 a_S=\nabla^S_{\dot\gamma_S}\dot\gamma_S
     =\mathcal K_{\pi_S}(\dot\Gamma,\dot\Gamma).
\]
In coordinates \(x^B\) on \(N\) and \(y^a\) on \(N_S\), the
defect is quadratic in velocity, with components
\begin{equation}
 (\mathcal K_{\pi_S})^a{}_{BC}
 =\partial_B\partial_C\pi_S^a
  +(\Gamma^S)^a{}_{bc}\,
        \partial_B\pi_S^b\partial_C\pi_S^c
  -\partial_D\pi_S^a(\Gamma^{\rm tot})^D{}_{BC}.
 \label{vii:eq:observador-defecto-coordenadas}
\end{equation}
\end{proposition}

\begin{proof}
The chain rule gives
\(\ddot\gamma_S^a=\partial_B\pi_S^a\ddot\Gamma^B
 +\partial_B\partial_C\pi_S^a\dot\Gamma^B\dot\Gamma^C\).
Adding the contribution of \(\nabla^S\) and subtracting the image of
the total acceleration cancels the terms containing \(\ddot\Gamma\).
The remaining terms are
\eqref{vii:eq:observador-defecto-coordenadas} contracted with
\(\dot\Gamma^B\dot\Gamma^C\). When the total acceleration vanishes,
the stated equality follows.
\end{proof}

The geometric response factors through \(b\) when the connection,
the projection, and the kinematic datum used in the evaluation are
determined by \(\beta=b(x)\), or when those kinematic data are
held fixed in the comparison. Under these conditions,
\eqref{vii:eq:observador-defecto-coordenadas} gives the same result
for any representatives of a fibre;
Theorem~\ref{vii:thm:observador-factorizacion} constructs its unique
factor on \(B_{\rm ef}\). If the kinematic datum varies, it
forms part of the domain of the response being compared.

The self-inertial formulation describes a geodesic trajectory of the
joint state and an effective acceleration of its subsystem due
to the projection defect. The gravitational realization subsequently identifies
that connection, frame, and response with their physical
observables. Observer covariance, memory inscription, and
factorization through the boundary thus retain complementary functions:
the first determines the readout class, the second retains the
genealogy, and the third specifies the relational dependence of the
response.
